% META: {"id": "TSPE-SUITLIM-CO-033", "chapitre": "TSPE-SUITES-LIMITES", "type_objet": "corrige", "exercice_ref": "TSPE-SUITLIM-EX-033", "capacites_codes": ["C5"], "sources_inspiration": [], "status": "approved"}
% BEGIN-VERIFY
% from sympy import *
% n = symbols('n', integer=True, nonnegative=True)
% a = Rational(1, 9)
% # Q2 : borne binomiale d'ordre 2, valable pour n >= 2
% for k in range(2, 30):
%     assert (1 + a)**k >= Rational(k*(k-1), 2) * a**2
% # Q3 : majoration de u_n = n*q^n, valable pour n >= 2
% q = Rational(9, 10)
% for k in range(2, 30):
%     u_k = k * q**k
%     assert u_k <= Rational(162, k - 1)
% assert limit(n * q**n, n, oo) == 0
% assert limit(Rational(162, 1) / (n - 1), n, oo) == 0
% END-VERIFY
\begin{corrige}{TSPE-SUITLIM-EX-033}

\textbf{Q1.} $u_1=0{,}9$, $u_5 = 5 \times 0{,}9^5 \approx 2{,}95$, $u_{10} = 10 \times 0{,}9^{10} \approx 3{,}49$.

\commentaireMarge{La formule du binome : $(1+a)^n = \displaystyle\sum_{k=0}^{n}\binom{n}{k}a^k$. Pour $a>0$, chaque terme est positif : on peut en garder un seul comme minorant.}

\textbf{Q2.} Par la formule du binome, pour $a = \dfrac{1}{9} > 0$ et $n \geq 2$ :
\[
  (1+a)^n = \sum_{k=0}^{n}\binom{n}{k}a^k \geq \binom{n}{2}a^2 = \frac{n(n-1)}{2}a^2,
\]
car tous les termes de la somme sont positifs. D'ou $\left(\dfrac{1}{q}\right)^n \geq \dfrac{n(n-1)}{2}a^2$.

\textbf{Q3.} En inversant l'inegalite precedente (tous les termes sont strictement positifs) :
\[
  q^n \leq \frac{2}{n(n-1)a^2} = \frac{2 \times 81}{n(n-1)} = \frac{162}{n(n-1)}.
\]
Donc $0 \leq u_n = n \times q^n \leq \dfrac{162\,n}{n(n-1)} = \dfrac{162}{n-1}$ pour $n \geq 2$.

Or $\lim_{n \to +\infty} \dfrac{162}{n-1} = 0$. Par le theoreme des gendarmes, $\boxed{\lim_{n \to +\infty} u_n = 0}$.

\end{corrige}
